\documentclass[11pt,a4paper]{article}

\usepackage[T1]{fontenc}
\usepackage[utf8]{inputenc}
\usepackage{lmodern}
\usepackage{microtype}
\usepackage{geometry}
\usepackage{graphicx}
\usepackage[table]{xcolor}
\usepackage{array}
\usepackage{tabularx}
\usepackage{longtable}
\usepackage{booktabs}
\usepackage{enumitem}
\usepackage{tikz}
\usepackage{tcolorbox}
\tcbuselibrary{skins,breakable}
\usepackage{titlesec}
\usepackage{fancyhdr}
\usepackage{hyperref}
\usepackage{float}
\newcommand{\Needspace}[1]{}

\geometry{margin=2.2cm}
\setlength{\parindent}{0pt}
\setlength{\parskip}{0.75em}
\setlength{\headheight}{24pt}
\emergencystretch=2em
\hfuzz=6pt

\definecolor{ForestGreen}{HTML}{2E6F40}
\definecolor{DeepGreen}{HTML}{1F4D2B}
\definecolor{SoftGreen}{HTML}{EAF4EA}
\definecolor{WarmSand}{HTML}{F4EFE6}
\definecolor{AccentGold}{HTML}{C89B3C}
\definecolor{TextGray}{HTML}{4A4A4A}

\hypersetup{
    colorlinks=true,
    linkcolor=ForestGreen,
    urlcolor=ForestGreen,
    pdftitle={Informe de especies nativas},
    pdfauthor={}
}

\pagestyle{fancy}
\fancyhf{}
\lhead{\textcolor{TextGray}{Informe de especies nativas}}
\rhead{\textcolor{TextGray}{Bosque Miyawaki}}
\cfoot{\textcolor{TextGray}{\thepage}}
\renewcommand{\headrulewidth}{0.4pt}
\renewcommand{\headrule}{\hbox to\headwidth{\color{ForestGreen}\leaders\hrule height \headrulewidth\hfill}}

\titleformat{\section}
  {\Large\bfseries\color{DeepGreen}}
  {\thesection}{0.75em}{}
\titleformat{\subsection}
  {\large\bfseries\color{ForestGreen}}
  {\thesubsection}{0.75em}{}

% ---------------------------------------------------------------------------
% Iconos vectoriales reutilizables (sol = exposición, gota = riego)
% ---------------------------------------------------------------------------
\newcommand{\sunicon}{%
  \begin{tikzpicture}[baseline=-0.55ex,scale=0.16]
    \draw[AccentGold,line width=1.6pt] (0,0) circle (0.9);
    \foreach \a in {0,45,90,135,180,225,270,315}{
      \draw[AccentGold,line width=1.6pt] (\a:1.25) -- (\a:1.75);
    }
  \end{tikzpicture}%
}
\newcommand{\dropicon}{%
  \begin{tikzpicture}[baseline=-0.55ex,scale=0.16]
    \fill[ForestGreen] (0,1.6) .. controls (-1.15,0.1) and (-1.15,-0.85) ..
      (0,-1.6) .. controls (1.15,-0.85) and (1.15,0.1) .. (0,1.6) -- cycle;
  \end{tikzpicture}%
}

% ---------------------------------------------------------------------------
% \speciescard: ficha visual de una especie.
% #1 título breve  #2 ruta de imagen  #3 nombre común  #4 nombre científico
% #5 estrato / altura  #6 exposición solar  #7 demanda de riego
% #8 función principal en el espacio
%
% NOTA PARA QUIEN EDITE ESTE DOCUMENTO:
% Para agregar una especie nueva, duplica un bloque \subsection{...}
% completo (llamada a \speciescard + párrafo + itemize) y reemplaza la
% ruta de la imagen, el nombre común, el nombre científico y los datos
% de los 8 campos principales.
% ---------------------------------------------------------------------------
\newcommand{\speciescard}[8]{%
\begin{tcolorbox}[
    breakable,
    colback=white,
    colframe=ForestGreen,
    boxrule=0.8pt,
    arc=2mm,
    left=1mm,
    right=1mm,
    top=1mm,
    bottom=1mm,
    title={\parbox[t]{0.95\linewidth}{\raggedright\textbf{#1}}},
    fonttitle=\bfseries,
    coltitle=white
]
\begin{minipage}[t]{0.34\textwidth}
\vspace{0pt}
    \centering
    \IfFileExists{#2}{%
        \includegraphics[width=\linewidth,height=5.5cm,keepaspectratio]{#2}%
    }{%
        \begin{tikzpicture}
            \fill[SoftGreen] (0,0) rectangle (5.5,5.5);
            \draw[ForestGreen, thick] (0.4,0.4) rectangle (5.1,5.1);
            \node[align=center, text=DeepGreen] at (2.75,2.75) {Imagen de la especie\\\small (#2)};
        \end{tikzpicture}%
    }
\end{minipage}
\hfill
\begin{minipage}[t]{0.62\textwidth}
\vspace{0pt}
\renewcommand{\arraystretch}{1.25}
\begin{tabularx}{\linewidth}{>{\bfseries\color{ForestGreen}}l >{\raggedright\arraybackslash}X}
Nombre común: & #3 \\
Nombre científico: & \textit{#4} \\
Estrato / altura: & #5 \\
\sunicon\ Exposición: & #6 \\
\dropicon\ Riego: & #7 \\
Función en el espacio: & #8
\end{tabularx}
\end{minipage}
\end{tcolorbox}
}

\begin{document}

\hypersetup{pageanchor=false}
\begin{titlepage}
    \begin{tikzpicture}[remember picture,overlay]
        \fill[WarmSand] (current page.south west) rectangle (current page.north east);
        \fill[ForestGreen] (current page.north west) rectangle ([xshift=6.2cm]current page.south west);
        \fill[SoftGreen] ([xshift=4.8cm,yshift=-3cm]current page.north west) circle (3.8cm);
        \fill[AccentGold, opacity=0.22] ([xshift=12cm,yshift=-13cm]current page.north west) circle (4.6cm);
        \draw[white, line width=1.4pt, opacity=0.85] ([xshift=1.3cm,yshift=-3cm]current page.north west) -- ([xshift=4.8cm,yshift=-3cm]current page.north west);
        \draw[white, line width=1.4pt, opacity=0.85] ([xshift=1.3cm,yshift=-4cm]current page.north west) -- ([xshift=4.0cm,yshift=-4cm]current page.north west);
    \end{tikzpicture}

    \vspace*{2.2cm}
    {\color{black}\large\bfseries Informe técnico}\par
    \vspace{0.8cm}
    {\color{black}\Huge\bfseries Especies nativas para plantación}\par
    \vspace{0.8cm}
    {\color{black}\Large Selección visual y fichas de especies para un espacio de restauración}\par

    \vspace{2.5cm}
    \begin{tcolorbox}[
        width=0.72\textwidth,
        colback=white,
        colframe=ForestGreen,
        arc=3mm,
        boxrule=0.9pt,
        left=4mm,
        right=4mm,
        top=3mm,
        bottom=3mm
    ]
    \textbf{Objetivo del documento}\par
    Presentar las especies nativas y endémicas que se utilizarán en la plantación de Tierra Amigos. El informe permite comparar las especies, conocer sus requerimientos y funciones ecológicas.
    \end{tcolorbox}

    \vfill
    {\color{white}\large\textbf{Proyecto:} Bosque Miyawaki}\par
\end{titlepage}
\hypersetup{pageanchor=true}

\tableofcontents
\newpage

\section{Introducción}
Se seleccionaron plantas nativas y endémicas de bajo requerimiento hídrico, capaces de adaptarse y prosperar en el contexto de escasez hídrica y sequía prolongada que afecta a la Región Metropolitana. Estas especies cumplen un rol ecológico relevante, ya que funcionan como refugio, hábitat y fuente de alimento para la fauna local, especialmente para polinizadores, como abejas nativas y mariposas, además de diversas aves.

Asimismo, algunas de estas especies poseen significados culturales asociados a usos medicinales, comestibles o simbólicos, lo que favorece la educación ambiental y fortalece el sentido de pertenencia de la comunidad con el jardín.

La selección se fundamenta en un monitoreo científico desarrollado durante varios meses, en el cual se evaluó positivamente la sobrevivencia de las especies, con valores superiores al 40\%, junto con su productividad y tolerancia al estrés en condiciones urbanas. Además, estas plantas presentan alta resistencia a plagas y una respuesta favorable a manejos agronómicos básicos, como podas de limpieza o de formación.

\section{Cómo usar esta guía}
Antes de revisar las fichas, esta sección explica brevemente los símbolos y términos que se repiten a lo largo del documento.

\begin{tcolorbox}[
    colback=SoftGreen,
    colframe=ForestGreen,
    arc=2mm,
    boxrule=0.6pt,
    left=4mm,right=4mm,top=3mm,bottom=3mm
]
\begin{tabularx}{\linewidth}{c X}
\sunicon & \textbf{Exposición solar}: cantidad de luz directa que necesita la especie (pleno sol, semisombra, etc.). \\[0.4em]
\dropicon & \textbf{Riego}: demanda hídrica orientativa durante la temporada de establecimiento (octubre--abril). \\
\end{tabularx}
\end{tcolorbox}

\section{Resumen comparativo de especies}
La siguiente tabla permite comparar de un vistazo las doce especies incluidas en este informe antes de revisar cada ficha en detalle.

\footnotesize
\renewcommand{\arraystretch}{1.35}
\setlength{\tabcolsep}{4pt}
\rowcolors{2}{SoftGreen}{white}
\begin{longtable}{>{\bfseries}p{2.7cm} p{2.1cm} p{1.7cm} p{2.1cm} p{1.8cm} p{2.6cm}}
\rowcolor{ForestGreen}
\textcolor{white}{\bfseries Especie} & \textcolor{white}{\bfseries Estrato} & \textcolor{white}{\bfseries Altura} & \textcolor{white}{\bfseries Exposición} & \textcolor{white}{\bfseries Riego} & \textcolor{white}{\bfseries Floración} \\
\endfirsthead
\rowcolor{ForestGreen}
\textcolor{white}{\bfseries Especie} & \textcolor{white}{\bfseries Estrato} & \textcolor{white}{\bfseries Altura} & \textcolor{white}{\bfseries Exposición} & \textcolor{white}{\bfseries Riego} & \textcolor{white}{\bfseries Floración} \\
\endhead
Corontillo & Arbusto / arbolito & Variable & Pleno sol & Moderado & Nov.--feb. \\
Huingán & Arbusto alto & Variable & Pleno sol & Moderado & Oct.--dic. \\
Quebracho & Arbusto / arbolito & Variable & Semisombra & Moderado & Primavera / feb. \\
Pingo Pingo & Arbusto & Variable & Pleno sol & Bajo & -- \\
Malva de cerro & Subarbusto & Hasta 1,4 m & Pleno sol & Bajo & Primavera--otoño \\
Maravilla de campo & Arbusto & Hasta 2 m & Pleno sol & Bajo & Primavera \\
Ruby & Hierba / subarbusto & Variable & Pleno sol & Bajo (goteo) & Fin invierno--otoño \\
Mirabilis ovata & Hierba perenne & Baja & Pleno sol & Muy bajo & -- \\
Pata de guanaco & Herb. suculenta & Baja & Pleno sol & Muy bajo & Primavera--verano \\
Poquil & Herbácea & Baja & Pleno sol & Bajo & Tras lluvias \\
\end{longtable}
\normalsize

\section{Especies arbustivas y arborescentes}

\Needspace{7cm}
\subsection{Corontillo}

\speciescard
{Corontillo}
{imagenes/corontillo.jpg}
{Corontillo}
{Escallonia pulverulenta}
{Arbusto o pequeño árbol}
{Pleno sol}
{Moderado ($\sim$8 L/semana)}
{Sombra, estructura vertical y refugio para fauna}

El corontillo es un arbusto o pequeño árbol endémico de Chile, caracterizado por su rápido crecimiento y su capacidad para desarrollarse en ambientes áridos y asoleados. En la Región Metropolitana se distribuye naturalmente en laderas secas, márgenes de bosques esclerófilos y orillas de caminos, por lo que presenta una buena compatibilidad con jardines de bajo requerimiento hídrico.

\begin{itemize}[leftmargin=1.2cm]
    \item \textbf{Ciclo de floración:} florece principalmente entre noviembre y febrero, formando densas espigas de flores blancas que atraen polinizadores, especialmente abejas y dípteros.
    \item \textbf{Fructificación:} produce frutos en forma de cápsula ovoidal, que contienen semillas pequeñas y alargadas.
    \item \textbf{Suelo:} se recomienda plantarlo en suelos pedregosos o rocosos, ya que estos favorecen la maduración de las cápsulas y la producción de semillas viables.
    \item \textbf{Exposición:} requiere pleno sol y se adapta bien a laderas áridas y sectores abiertos con alta radiación.
    \item \textbf{Crecimiento:} es una especie de rápido crecimiento que puede alcanzar hasta 5 metros de altura, aportando estructura vertical y sombra al espacio.
    \item \textbf{Compatibilidad con el sitio:} es adecuada para espacios urbanos secos y soleados. Sin embargo, se debe evitar su plantación en suelos excesivamente fértiles o blandos, ya que puede desarrollar abundante follaje, pero reducir su capacidad reproductiva.
\end{itemize}

\newpage

\Needspace{7cm}
\subsection{Huingán}

\speciescard
{Huingán}
{imagenes/huingan.jpg}
{Huingán}
{Schinus polygamus}
{Arbusto alto o arbolito}
{Pleno sol}
{Moderado ($\sim$8 L/semana)}
{Sombra, estructura vertical y refugio para fauna}

El huingán es un arbusto alto, ampliamente adaptado a condiciones de sequía, alta exposición solar y suelos pobres. En la Región Metropolitana es frecuente en valles interiores secos y forma parte característica de la estepa de espino, donde convive con especies como el peumo, el quillay y el espino.

\begin{itemize}[leftmargin=1.2cm]
    \item \textbf{Ciclo de floración:} florece entre la primavera media y el inicio del verano, aproximadamente entre octubre y diciembre, produciendo pequeñas flores amarillentas.
    \item \textbf{Fructificación:} sus frutos maduran hacia finales del verano, entre enero y marzo. Corresponden a drupas violáceas oscuras o negras brillantes, que sirven de alimento para aves locales.
    \item \textbf{Multiplicación:} se multiplica con facilidad por semillas y presenta un crecimiento relativamente rápido, lo que favorece su establecimiento en proyectos de restauración o jardinería nativa.
    \item \textbf{Crecimiento:} suele presentar un hábito arbustivo, alcanzando hasta 3 metros de altura. Sus ramas son intrincadas y, en algunos casos, pueden terminar en espinas.
    \item \textbf{Suelo y clima:} prefiere suelos pobres, áridos, degradados y bien expuestos al sol. Tolera condiciones exigentes propias de la zona central, incluyendo heladas y temperaturas extremas.
    \item \textbf{Compatibilidad con el sitio:} es una especie adecuada para espacios urbanos secos y soleados, especialmente en sectores donde se busca incorporar vegetación resistente, de bajo mantenimiento y con valor ecológico para la fauna local.
\end{itemize}

\newpage

\Needspace{7cm}
\subsection{Quebracho}

\speciescard
{Quebracho}
{imagenes/quebracho.jpg}
{Quebracho}
{Senna candolleana}
{Arbusto o arbolito}
{Semisombra}
{Moderado ($\sim$8 L/semana)}
{Sombra, estructura vertical, recuperación de suelo y refugio para fauna}

El quebracho es una especie endémica de Chile, presente desde la Región de Coquimbo hasta la Región Metropolitana. Habita principalmente en valles secos y laderas expuestas, donde cumple un rol relevante en la estabilización y recuperación de ambientes degradados. Bajo el nombre común de quebracho pueden encontrarse especies afines del género \textit{Senna}.

\begin{itemize}[leftmargin=1.2cm]
    \item \textbf{Ciclo de floración:} su floración principal ocurre durante la primavera, produciendo flores muy vistosas de color amarillo anaranjado. Como referencia, la especie afín \textit{Senna stipulacea} florece específicamente en febrero.
    \item \textbf{Fructificación:} produce frutos en forma de legumbre cilíndrica, de color pardo al madurar.
    \item \textbf{Crecimiento:} presenta un hábito arbustivo globoso y puede alcanzar hasta 6 metros de altura, aportando estructura vertical y cobertura al espacio.
    \item \textbf{Siembra:} se propaga por semillas, las cuales requieren estratificación previa para favorecer su germinación.
    \item \textbf{Beneficio edáfico:} posee la capacidad de fijar nitrógeno en el suelo mediante simbiosis radicular, lo que contribuye a mejorar la fertilidad del sitio y favorece el establecimiento de otras especies vegetales.
    \item \textbf{Sensibilidad:} es una especie sensible a las heladas, por lo que se recomienda ubicarla en sectores protegidos, especialmente durante sus primeras etapas de establecimiento.
    \item \textbf{Compatibilidad con el sitio:} es adecuada para procesos de restauración de laderas degradadas y espacios urbanos secos. Para favorecer su establecimiento, se recomienda elegir un lugar a semisombra, especialmente al momento del trasplante.
    \item \textbf{Manejo agronómico recomendado:} en la Región Metropolitana se recomienda aplicar riego activo entre octubre y abril, con un volumen aproximado de 8 litros semanales por planta. El riego debe aplicarse lejos del cuello del tronco para evitar problemas asociados a hongos. Durante los meses de mayor calor, especialmente enero y febrero, se sugiere dividir el riego en dos aplicaciones semanales.
\end{itemize}

\newpage

\section{Especies xerófitas y herbáceas}

\Needspace{7cm}
\subsection{Pingo Pingo}

\speciescard
{Pingo Pingo}
{imagenes/pingo-pingo.jpg}
{Pingo Pingo}
{Ephedra chilensis}
{Arbusto nativo postrado o erecto}
{Pleno sol}
{Bajo ($\sim$4 L/semana)}
{Cobertura de suelo, atracción de polinizadores y protección del sustrato}

El pingo pingo es un arbusto nativo de hábitos postrados o erectos, altamente adaptado a condiciones de aridez, sequía y suelos pedregosos. En la Región Metropolitana se distribuye naturalmente en sectores como Río Clarillo y el Cajón del Maipo, donde crece asociado a ambientes secos, erosionados y de alta exposición solar.

\begin{itemize}[leftmargin=1.2cm]
    \item \textbf{Crecimiento:} es un arbusto ramoso que puede alcanzar entre 30 y 120 cm de altura. Presenta una alta productividad de biomasa cuando se desarrolla en suelos pedregosos, independientemente del sistema de riego utilizado.
    \item \textbf{Ciclo biológico:} al ser una gimnosperma, posee hojas reducidas a pequeñas escamas, lo que le permite disminuir la pérdida de agua y adaptarse mejor a condiciones de sequía.
    \item \textbf{Fenología:} produce estructuras reproductivas llamadas estróbilos. Los masculinos suelen disponerse en tríadas, mientras que los femeninos contienen entre 5 y 7 semillas oval-lanceoladas.
    \item \textbf{Suelo y plantación:} se recomienda plantarlo en sustratos pedregosos o erosionados, ya que estas condiciones favorecen su desarrollo, floración y fructificación. Es importante evitar suelos excesivamente enriquecidos con materia orgánica, porque pueden estimular el crecimiento vegetativo de tallos e inhibir la floración.
    \item \textbf{Multiplicación:} puede reproducirse mediante semillas o por fraccionamiento de rizomas, lo que facilita su propagación en proyectos de restauración o jardinería nativa.
    \item \textbf{Riego:} requiere un bajo aporte hídrico, con un volumen aproximado de 4 litros semanales por planta durante la etapa de establecimiento.
\end{itemize}

\newpage

\Needspace{7cm}
\subsection{Malva de cerro}

\speciescard
{Malva de cerro}
{imagenes/malva-de-cerro.jpg}
{Malva de cerro}
{Sphaeralcea obtusiloba}
{Subarbusto perenne}
{Pleno sol}
{Bajo ($\sim$4 L/semana)}
{Cobertura de suelo, atracción de polinizadores y protección del sustrato}

La malva de cerro es un subarbusto endémico de Chile, perteneciente a la familia de las malváceas. Su distribución natural se restringe principalmente a las regiones de Valparaíso, Metropolitana y O'Higgins, por lo que presenta una alta pertinencia ecológica para proyectos de jardinería nativa en la zona central. Es una especie siempreverde, resistente y muy valorada por su prolongada floración y su capacidad para atraer abejas nativas.

\begin{itemize}[leftmargin=1.2cm]
    \item \textbf{Ciclo de floración:} presenta una floración prolongada, que puede extenderse desde el invierno tardío hasta el final de la primavera e incluso continuar hacia el otoño. Sus flores son vistosas, de tonalidades violeta o púrpura pálido.
    \item \textbf{Crecimiento:} es un subarbusto siempreverde que puede alcanzar hasta 1,4 metros de altura, con un diámetro de copa similar. Sus hojas presentan un tono verde grisáceo debido a la presencia de pelos, los cuales ayudan a reflejar la radiación solar.
    \item \textbf{Suelo y plantación:} se adapta a una amplia variedad de texturas y calidades de suelo, siempre que cuente con buen drenaje. Esta plasticidad agroecológica permite su establecimiento en distintos sectores de la Región Metropolitana.
    \item \textbf{Manejo:} responde favorablemente a una poda anual de renovación, la cual estimula la aparición de brotes florales vigorosos y evita que las ramas se vuelvan excesivamente leñosas.
    \item \textbf{Riego:} presenta una demanda hídrica baja, con un volumen aproximado de 4 litros semanales por planta durante la etapa de establecimiento.
    \item \textbf{Compatibilidad con el sitio:} requiere pleno sol y se adapta bien a condiciones urbanas secas, siempre que el suelo tenga buen drenaje. Es una especie especialmente adecuada para aportar cobertura, color estacional y recursos florales para polinizadores nativos.
\end{itemize}

\newpage

\Needspace{7cm}
\subsection{Maravilla de campo o Incienso}

\speciescard
{Maravilla de campo}
{imagenes/maravilla-de-campo.jpg}
{Maravilla de campo o Incienso}
{Flourensia thurifera}
{Arbusto hasta 2 m}
{Pleno sol}
{Bajo (4 a 8 L/semana)}
{Cobertura de suelo, atracción de polinizadores y protección del sustrato}

La maravilla de campo, también conocida como incienso, es un arbusto endémico de Chile, presente desde la Región de Coquimbo hasta La Araucanía. En la Región Metropolitana es común en laderas secas y márgenes de caminos, donde destaca por su alta rusticidad y capacidad para establecerse en ambientes degradados. Es una especie pionera, ideal para la recuperación de suelos erosionados, especialmente en laderas con alta exposición solar.

\begin{itemize}[leftmargin=1.2cm]
    \item \textbf{Ciclo de floración:} su floración principal ocurre durante el invierno y la primavera, aunque puede extenderse hasta el otoño según las condiciones del sitio. Produce capítulos florales amarillos, vistosos y de gran tamaño, que pueden alcanzar hasta 8 cm de diámetro.
    \item \textbf{Ciclo biológico:} una característica particular de esta especie es que sus hojas se marchitan durante el verano, pero permanecen adheridas a las ramas. Esta adaptación le permite enfrentar períodos de alta sequedad ambiental.
    \item \textbf{Crecimiento:} es un arbusto erecto, muy ramificado, que puede alcanzar hasta 2 metros de altura. Sus hojas exudan una resina aromática que actúa como barrera impermeable, ayudando a reducir la pérdida de humedad.
    \item \textbf{Suelo y plantación:} se adapta bien a suelos pobres, secos, erosionados o degradados, por lo que resulta adecuada para laderas, bordes de caminos y sectores con baja disponibilidad hídrica.
    \item \textbf{Usos:} es valorada en apicultura por su aporte floral para polinizadores. Además, posee usos en la medicina tradicional y su resina ha sido utilizada históricamente como incienso.
    \item \textbf{Riego:} requiere bajo aporte hídrico, con un volumen aproximado de entre 4 y 8 litros semanales por planta, dependiendo de su tamaño y etapa de establecimiento.
    \item \textbf{Compatibilidad con el sitio:} presenta una alta plasticidad agroecológica, adaptándose a distintas condiciones de la Región Metropolitana. Su carácter pionero la hace especialmente útil para colonizar suelos erosionados o degradados, aportar cobertura vegetal y favorecer la llegada de polinizadores.
\end{itemize}

\newpage

\Needspace{7cm}
\subsection{Ruby}
\speciescard{Ruby}{imagenes/ruby.jpg}{Ruby}{Alternanthera porrigens}{Hierba perenne o subarbusto}{Pleno sol}{Bajo, riego localizado}{Cobertura de suelo, atracción de polinizadores y protección del sustrato}{}
Es una hierba perenne o subarbusto muy valorado por el color intenso de sus inflorescencias.
\begin{itemize}[leftmargin=1.2cm]
     \item \textbf{Ciclo biológico:} produce pequeños capítulos florales de color púrpura a rojo intenso (tono ``ruby''), muy vistosos, y mantiene su follaje persistente durante todo el año.
    \item \textbf{Ciclo de floración:} se extiende desde el final del invierno hasta el otoño.
    \item \textbf{Condiciones de plantación:} su productividad se optimiza en suelos pedregosos; es muy sensible al exceso de humedad.
    \item \textbf{Exposición:} requiere pleno sol y es ideal para ambientes costeros.
    \item \textbf{Riego:} se recomienda riego localizado por goteo, evitando mojar el follaje por aspersión.
\end{itemize}

\Needspace{7cm}
\subsection{Mirabilis ovata}
\speciescard{Mirabilis ovata}{imagenes/mirabilis-ovata.jpg}{Mirabilis ovata}{Mirabilis ovata}{Hierba perenne xerófila}{Pleno sol}{Muy bajo}{Cobertura de suelo y aporte de biodiversidad floral}{}
Es una hierba perenne xerófila nativa, adaptada a condiciones de aridez extrema.
\begin{itemize}[leftmargin=1.2cm]
    \item \textbf{Estrategia de supervivencia:} posee una raíz tuberosa napiforme que reserva agua y energía.
    \item \textbf{Ciclo reproductivo:} puede autopolinizarse enrollando los filamentos estaminales hacia el estigma.
    \item \textbf{Hábitat de plantación:} se adapta a laderas áridas y afloramientos graníticos desde el nivel del mar hasta los 3.200 m s.\,n.\,m.
\end{itemize}

\Needspace{7cm}
\subsection{Pata de guanaco}
\speciescard{Pata de guanaco}{imagenes/pata-de-guanaco.jpg}{Pata de guanaco}{Cistanthe grandiflora}{Herbácea suculenta perenne}{Pleno sol}{Muy bajo}{Cobertura de suelo y aporte de biodiversidad floral}{}
Es una herbácea suculenta perenne, extremadamente resistente y de gran valor ornamental.
\begin{itemize}[leftmargin=1.2cm]
    \item \textbf{Ciclo de floración:} produce flores magenta brillante principalmente durante primavera y verano.
    \item \textbf{Multiplicación:} se propaga con gran facilidad mediante semillas y esquejes.
    \item \textbf{Resistencia:} puede sobrevivir más de dos meses sin riego en pleno verano; tolera heladas de hasta -5\,\textdegree C, pero no nieve.
\end{itemize}

\Needspace{7cm}
\subsection{Poquil}
\speciescard{Poquil}{imagenes/poquil.jpg}{Poquil}{Helenium aromaticum / H. glaucum}{Herbácea del secano semiárido}{Pleno sol}{Bajo}{Cobertura de suelo y aporte de biodiversidad floral}{}
Este complejo de especies herbáceas es típico del secano semiárido.
\begin{itemize}[leftmargin=1.2cm]
    \item \textbf{Ciclo de floración:} florece con vigor inmediatamente después de las escasas lluvias de invierno.
    \item \textbf{Condiciones de plantación:} requiere suelos con excelente drenaje y baja humedad edáfica.
    \item \textbf{Características:} es muy resistente a la degradación del suelo y posee glándulas con aceites esenciales de olor intenso y agradable.
\end{itemize}

\section{Protocolo general de riego}
Todas las especies de este informe comparten un mismo calendario de riego activo, concentrado entre octubre y abril, con doble frecuencia semanal durante enero y febrero. El volumen de agua varía según el estrato de la especie.

\begin{center}
\renewcommand{\arraystretch}{1.4}
\begin{tabularx}{\linewidth}{>{\bfseries\color{ForestGreen}}p{4.6cm} X X}
\toprule
 & \textbf{Estrato alto} (árboles, arbustos y arborescentes) & \textbf{Estrato bajo} (herbáceas y subarbustos) \\
\midrule
Periodo activo & Octubre a abril & Octubre a abril \\
Frecuencia en verano & Dos veces por semana (enero--febrero) & Dos veces por semana (enero--febrero) \\
Volumen semanal & $\sim$8 litros por planta & $\sim$4 litros por planta \\
Método & Riego localizado en la zona de raíces & Aplicar lejos del cuello de la planta para evitar hongos radiculares \\
\bottomrule
\end{tabularx}
\end{center}

\section{Complementación entre especies}
La plantación cercana de estas especies puede organizarse como una asociación coherente, en la que cada grupo cumple una función ecológica complementaria. El objetivo es aprovechar requerimientos similares de agua y suelo, al mismo tiempo que se refuerzan procesos como la polinización, la protección del suelo y la estabilidad del sistema.

\subsection{Estratificación vertical}
La combinación de árboles como Huingán, Corontillo y Quebracho con arbustos como Malva de cerro, Maravilla de campo y Pingo Pingo, más herbáceas como Ruby, Pata de guanaco y Poquil, permite construir una estratificación vertical diversa.
\begin{itemize}[leftmargin=1.2cm]
    \item \textbf{Refugio para fauna:} la superposición de capas vegetales ofrece sitios de anidación, alimentación y resguardo para insectos, aves y reptiles.
    \item \textbf{Microclima:} la diversidad de alturas ayuda a moderar la temperatura del suelo y a generar un ambiente más fresco y protegido.
\end{itemize}

\subsection{Sinergia nutricional y estabilización}
Las especies seleccionadas no solo conviven, sino que también se potencian entre sí mediante procesos de mejora del suelo y control de erosión.
\begin{itemize}[leftmargin=1.2cm]
    \item \textbf{Fijación de nitrógeno:} el Quebracho, como especie del género \textit{Senna}, contribuye a aumentar la fertilidad del suelo.
    \item \textbf{Estabilidad del sustrato:} especies arbustivas como Corontillo y Malva de cerro ayudan a fijar laderas degradadas y protegen a las herbáceas más pequeñas.
\end{itemize}

\subsection{Control biológico de plagas}
La diversidad vegetal favorece el equilibrio natural del sistema y reduce la dependencia de manejos externos.
\begin{itemize}[leftmargin=1.2cm]
    \item \textbf{Atracción de polinizadores:} Malva de cerro, Pata de guanaco y Poquil atraen abejas nativas, mariposas y sírfides.
    \item \textbf{Equilibrio ecológico:} una mayor biodiversidad de insectos benéficos ayuda a mantener bajo control a posibles plagas.
\end{itemize}

\subsection{Compatibilidad agroecológica}
La mayoría de estas especies se ubica en categorías CAE 1 y CAE 2, o bien en rangos de adaptación cercanos, lo que facilita su convivencia bajo un régimen de riego compartido.
\begin{itemize}[leftmargin=1.2cm]
    \item \textbf{Rusticidad compartida:} todas presentan bajos requerimientos hídricos y buena adaptación a condiciones de aridez o semiaridez.
    \item \textbf{Manejo común:} pueden mantenerse bajo un mismo programa de riego sin generar desequilibrios importantes entre especies.
\end{itemize}

En conjunto, la plantación mixta permite ahorrar agua y, al mismo tiempo, restituir servicios ecosistémicos clave como la polinización y la formación de suelo vivo.

\section{Especificaciones generales del espacio}
\begin{tabularx}{\linewidth}{>{\bfseries\color{ForestGreen}}p{4cm} X}
Tipo de plantación: & Alta diversidad de estratos y densidad intermedia a alta. \\
Objetivo principal: & Restauración ecológica y mejora paisajística. \\
Manejo inicial: & Riego de establecimiento, control de malezas y reposición de fallas. \\
Mantenimiento: & Monitoreo periódico de crecimiento, sanidad y competencia. \\
\end{tabularx}

\end{document}